% Figure: efficiency waterfall for the silicon solar-cell case study.
% Starting from incident 100%, successive bars subtract each loss term
% (below-gap transmission, above-gap thermalisation, detailed-balance
% voltage/recombination, real-cell losses) ending at the delivered
% module efficiency. Drawn with plain TikZ rectangles.
\begin{figure}[htbp]
\centering
\begin{tikzpicture}[font=\scriptsize]
% axis: y = efficiency (%) on a 0..100 scale, scaled by 0.06 cm per %
\def\sy{0.06}
% baseline and gridlines
\draw[enggray, thin] (0,0) -- (12.5,0);
\foreach \v in {0,20,40,60,80,100} {
  \draw[enggray, very thin] (-0.1,\v*\sy) -- (12.5,\v*\sy);
  \node[enggray, anchor=east, font=\tiny] at (-0.15,\v*\sy) {\v};
}
\node[enggray, anchor=south, rotate=90, font=\scriptsize]
  at (-0.9,50*\sy) {fraction of incident power (\%)};
% bars: each is (x-left, low, high, color, label-top, label-bottom)
% 1 incident 100
\filldraw[engblue]   (0.3,0)         rectangle (1.7,100*\sy);
\node[anchor=south, font=\tiny] at (1.0,100*\sy) {100};
\node[anchor=north, font=\tiny, align=center] at (1.0,0) {incident\\AM1.5G};
% 2 after below-gap loss: ~77 absorbable (drop of 23)
\filldraw[engamber]  (2.3,0)         rectangle (3.7,77*\sy);
\node[anchor=south, font=\tiny] at (3.0,77*\sy) {77};
\node[anchor=north, font=\tiny, align=center] at (3.0,0) {above-gap\\absorbed};
% 3 after thermalisation: ~49 left
\filldraw[engorange] (4.3,0)         rectangle (5.7,49*\sy);
\node[anchor=south, font=\tiny] at (5.0,49*\sy) {49};
\node[anchor=north, font=\tiny, align=center] at (5.0,0) {after\\thermal-\\isation};
% 4 detailed-balance ceiling ~33 (Shockley-Queisser)
\filldraw[engred]    (6.3,0)         rectangle (7.7,33*\sy);
\node[anchor=south, font=\tiny] at (7.0,33*\sy) {33};
\node[anchor=north, font=\tiny, align=center] at (7.0,0) {S--Q\\limit};
% 5 lab silicon record ~26
\filldraw[englime]   (8.3,0)         rectangle (9.7,26*\sy);
\node[anchor=south, font=\tiny] at (9.0,26*\sy) {26};
\node[anchor=north, font=\tiny, align=center] at (9.0,0) {lab Si\\record};
% 6 commercial module ~21
\filldraw[enggray]   (10.3,0)        rectangle (11.7,21*\sy);
\node[anchor=south, font=\tiny] at (11.0,21*\sy) {21};
\node[anchor=north, font=\tiny, align=center] at (11.0,0) {module};
\end{tikzpicture}
\caption[Silicon solar-cell efficiency waterfall]{Where the power goes in a
silicon photovoltaic cell, as a running fraction of the incident AM1.5G
irradiance. Below-gap photons pass through; the excess energy of above-gap
photons thermalises; the detailed-balance limit of unavoidable radiative
recombination and the voltage shortfall caps a single junction near
$33\,\%$; real-cell recombination, series resistance, and reflection bring a
record laboratory silicon cell to about $26\,\%$; module assembly, cell
spacing, and glass reflection bring a commercial panel near $21\,\%$. Each
step is a separate physical loss with a separate engineering remedy, and the
case study in this section carries them in order.}
\label{fig:vol04ch12:sc-cell-loss-waterfall}
\end{figure}
